\begin{Solbook}{1}
 \par \noindent  \begin {listEX} \item $(-27)+15-(-6)=(-27)+15+6= -6$. \item $(-24)\cdot 15+(-24)\cdot 35=(-24)\cdot (15+35)=(-24)\cdot 50=-1200$. \item $\left [(-36):4+2\cdot (-5)\right ]:(-1)^{2}+(-2)^{3}$\\ $=\left [-9+(-10)\right ]:1+(-8)$\\ $=(-19)+(-8)=-27$. \item $540-18:\left [(8-5)^{2}+(-2)^{3}\cdot 5^{0}\right ]$\\ $=540-18:\left [9+(-8)\cdot 1\right ]$\\ $=540-18:(9-8)$\\ $=540-18:1=522$. \end {listEX}  \hfill \qedEX 
\end{Solbook}
\begin{Solbook}{2}
 \par \noindent  \begin {listEX} \item $8x=-56$, suy ra $x=-56:8=-7$. \item $|3x+6|=15$ nên có hai trường hợp.\\ Trường hợp 1: $3x+6=15$, nên $3x=9$, $x=3$.\\ Trường hợp 2: $3x+6=-15$, nên $3x=-21$, $x=-7$.\\ Vậy $x=3$ hoặc $x=-7$. \item $4\left [2(x-3)+5\right ]-6=30$\\ $4(2x-6+5)-6=30$\\ $4(2x-1)-6=30$\\ $8x-4-6=30$\\ $8x-10=30$\\ $8x=40$\\ $x=5$. \item $4(x+3)-3(x-1)=2x+5$\\ $4x+12-3x+3=2x+5$\\ $x+15=2x+5$\\ $x=10$. \end {listEX}  \hfill \qedEX 
\end{Solbook}
\begin{Solbook}{3}
 \par \noindent  \begin {enumerate} \item Các số nguyên từ $-12$ đến $12$ chia hết cho $3$ là\\ $-12;\,-9;\,-6;\,-3;\,0;\,3;\,6;\,9;\,12$.\\ Trong các số này, số không chia hết cho $2$ là $-9;\,-3;\,3;\,9$.\\ Vậy $M=\{-9;\,-3;\,3;\,9\}$. \item Số vừa thuộc $M$ vừa thuộc $N$ là $3$ và $9$.\\ $N=\{0;\,3;\,6;\,9;\,12\}$.\\ Các số thuộc $N$ nhưng không thuộc $M$ là $0;\,6;\,12$. \end {enumerate}  \hfill \qedEX 
\end{Solbook}
\begin{Solbook}{4}
 \par \noindent  \begin {enumerate} \item Số tổ phải là ước chung của $36$ và $48$.\\ $36=2^{2}\cdot 3^{2}$, $48=2^{4}\cdot 3$, nên ước chung lớn nhất là $2^{2}\cdot 3=12$.\\ Các ước dương của $12$ là $1;\,2;\,3;\,4;\,6;\,12$.\\ Vì số tổ lớn hơn $1$, có $5$ cách: chia thành $2$, $3$, $4$, $6$ hoặc $12$ tổ. \item Mỗi tổ ít bạn nhất khi số tổ nhiều nhất, tức là $12$ tổ.\\ Mỗi tổ có $36:12=3$ bạn múa và $48:12=4$ bạn hát. \end {enumerate}  \hfill \qedEX 
\end{Solbook}
\begin{Solbook}{5}
 \par \noindent  Số sách là bội chung của $12$ và $18$.\\ $12=2^{2}\cdot 3$, $18=2\cdot 3^{2}$, nên bội chung nhỏ nhất là $2^{2}\cdot 3^{2}=36$.\\ Các bội của $36$ là $0;\,36;\,72;\,108;\,144;\,180;\,216;\ldots $\\ Số sách $n$ thỏa mãn $140<n<180$, nên $n=144$.\\ Thử lại: $144:12=12$, $144:18=8$.\\ Vậy thư viện quyên góp được $144$ cuốn sách.  \hfill \qedEX 
\end{Solbook}
\begin{Solbook}{6}
 \par \noindent  \begin {enumerate} \item Diện tích bồn hoa là $\dfrac {(10+6)\cdot 4}{2}=32$ (m$^{2}$). \item Diện tích sân là $16\cdot 10=160$ (m$^{2}$).\\ Diện tích lát gạch là $160-32=128$ (m$^{2}$). \item Số tiền lát gạch là $128\cdot 125\,000=16\,000\,000$ (đồng). \end {enumerate}  \hfill \qedEX 
\end{Solbook}
\begin{Solbook}{7}
 \par \noindent  $B=7^{4}-7^{2}=7^{2}\cdot (7^{2}-1)=49\cdot (49-1)=49\cdot 48$.\\ Vì $B=49\cdot 48$ nên $B$ chia hết cho $48$.  \hfill \qedEX 
\end{Solbook}
